%قبل از شروع مسئله ترجیحاً چیزی ننویسید. اما اگر خواستید بنویسید، در قالب لیستی (فقرات و شمارش) بنویسید.
\شروع‌مسئله[12]
% صورت اصلی مسئله را در این‌جا بنویسید.
\\
شما با شنگدباو بازی‌ای دارید می‌کنید! شنگدباو در ابتدا $n$ تا گوی با شماره‌های $1$ تا $n$ دارد. $m$ مرحله و در هر سری شما به او طنابی می‌دهید و از او می‌خواهید گوی $u$ و $v$ را به هم وصل کند (فرض کنید $u \neq v$). شنگدباو دوست دارد گوی‌هایش را رنگ کند و تنها دو رنگ آبی و قرمز را برای اینکار استفاده می‌کند. او به نظرش یک رنگ‌آمیزی زیبا است اگر هیچ طنابی نباشد که گوی‌های دو سرش رنگ برابر داشته باشند.

بنابراین هر سری که طنابی به او می‌دهید او رنگ تعدادی گوی را عوض می‌کند تا شرط زیبایی برقرار شود (در ابتدا همه گوی‌ها آبی‌اند) اگر هم نتوان چنین رنگ‌آمیزی‌ای انجام داد درخواست شما را رد می‌کند و طناب مورد نظر را پاره‌پاره می‌کند و سپس آتش می‌زند و خاکستر آن‌را در باغچه می‌ریزد! 
عملیات از بین بردن طناب‌ها شنگدباو را خسته نمی‌کند اما هر گویی که رنگ می‌کند و یا رنگ آن‌را عوض می‌کند ۱ واحد خسته می‌شود. 

روشی به شنگدباو ارائه دهید تا حداکثر
$O((n + m) log_2n)$
مقدار خسته شود.
%اگر مسئله چند قسمت داشت به این صورت عمل کنید:
%\شروع{زیرمسائل}
%\زیرمسئله[0] % مسئلۀ ۱آ
%\زیرمسئله[0] % مسئلۀ ۱ب
%\زیرمسئله[0] % مسئلۀ ۱پ
%\زیرمسئله[0] % مسئلۀ ۱ت
%\پایان{زیرمسائل}



%%%%%%%%%%%%%%%%%%%%%%%%%%%%%%%%%%%%%%%%%%%%%%%%%%%%%
% برخی دستورات مفید برای تمرین‌ها. %
%%%%%%%%%%%%%%%%%%%%%%%%%%%%%%%%%%%%%%%%%%%%%%%%%%%%%

% فرمول‌های نمونه:
%$f(x) = \Theta(x)$
%$T(x, y) \leq \log(x) + 2^y + \binom{2n}{n}$
% {\tt code\_function}


% لیست موارد بدون شماره‌گذاری:
%\شروع{فقرات}
%    \فقره مورد اول 
%        \شروع{فقرات}
%            \فقره مورد اول داخلی 
%                \شروع{فقرات}
%                    \فقره اولین مورد خیلی داخلی 
%                    \فقره دومین مورد خیلی داخلی 
%                \پایان{فقرات}
%            \فقره مورد دوم داخلی 
%        \پایان{فقرات}
%    \فقره مورد دوم 
%\پایان{فقرات}

% لیست موارد شماره‌گذاری شده:
%\شروع{شمارش}
%    \فقره یک مورد 
%    \فقره یک مورد دیگر 
%    \فقره باز هم یک مورد دیگر
%\پایان{شمارش}

% برای نوشتن یک فرمول چند خطی که یک شماره بخوره. برای حالت‌های دیگر به سایت زیر مراجعه کنید:
% https://www.overleaf.com/learn/latex/Aligning_equations_with_amsmath
%\begin{align} 
%    \begin{split}
%        (x+y)^3 &= (x+y)^2(x+y) \\
%                &= (x^2+2xy+y^2)(x+y) \\
%                &= (x^3+2x^2y+xy^2) + (x^2y+2xy^2+y^3) \\
%                &= x^3+3x^2y+3xy^2+y^3
%    \end{split}                                 
%\end{align}

% نوشتن ماتریس و آرایه دوبعدی
%\begin{align}
%    A = 
%    \begin{bmatrix}
%        A_{11} & A_{21} \\
%        A_{21} & A_{22}
%    \end{bmatrix}
%\end{align}

% اضافه کردن تصویر با شرح. برای اضافه کردن تصویر بدون شرح، فقط دستور درج‌تصویر را استفاده کنید.
%\شروع{شکل}[h!]
%\تنظیم‌ازوسط
%  \درج‌تصویر[پهنا=0.5\پهنای‌متن]{img.jpg}
%\شرح{شرح تصویر}
%\برچسب{تصویر:۱}
%\پایان{شکل}
%\پایان{زیرمسائل}

%کد زیر را ببینید.
%%\lstinputlisting{source_filename.py}
%\begin{latin}
%\begin{lstlisting}
%# Import the modules
%import sys
%import random
%
%ans = True
%
%while ans:
%question = raw_input("Ask the magic 8 ball a question: (press enter to quit) ")
%
%answers = random.randint(1,8)
%
%if question == "":
%sys.exit()
%
%elif answers == 1:
%print "It is certain"
%
%elif answers == 2:
%print "Outlook good"
%
%elif answers == 3:
%print "You may rely on it"
%
%elif answers == 4:
%print "Ask again later"
%
%elif answers == 5:
%print "Concentrate and ask again"
%
%elif answers == 6:
%print "Reply hazy, try again"
%
%elif answers == 7:
%print "My reply is no"
%
%elif answers == 8:
%print "My sources say no"
%\end{lstlisting}
%\end{latin}
